\section{Phase polynomials}\label{app:phase-polynomials}

This appendix collects the phase-polynomial facts used throughout the paper.
For $A\subseteq [n]$, write $x_A=\prod_{i\in A}x_i$, and regard the
binary inner product $\langle u,x\rangle$ as an element of
\(\{0,1\}\subseteq\Z_{2^k}\).

\begin{theorem}[Equivalent forms of phase polynomials]
    \label{thm:phase-polynomial-equivalence}
    Let $k\geq 1$.  A function $q:\F_2^n\to\Z_{2^k}$ can be written in
    the weighted form
    \begin{equation}
        \label{eq:weighted-phase-polynomial}
        q(x)=c_\varnothing+
        \sum_{\substack{\varnothing\neq A\subseteq[n]\\ |A|\leq k}}
        2^{|A|-1}c_Ax_A
        \pmod {2^k}
    \end{equation}
    if and only if it can be written as a sum of parity functions,
    \begin{equation}
        \label{eq:phase-polynomial-sum-parities}
        q(x)=a_0+
        \sum_{u\in\F_2^n\setminus\{0\}}a_u\langle u,x\rangle
        \pmod {2^k}.
    \end{equation}
    All coefficients may be taken in \(\Z_{2^k}\), although a coefficient
    $c_A$ with $A\neq\varnothing$ only matters modulo
    $2^{k-|A|+1}$.  In particular, for $k=3$, the weighted form is
    \[
        c_\varnothing+\sum_i c_i x_i
        +2\sum_{i<j}c_{ij}x_ix_j
        +4\sum_{i<j<\ell}c_{ij\ell}x_ix_jx_\ell
        \pmod 8.
    \]
\end{theorem}

\begin{proof}
    Repeatedly applying $a\oplus b=a+b-2ab$ gives, for every nonempty
    $B\subseteq[n]$,
    \begin{equation}
        \label{eq:parity-to-weighted}
        \bigoplus_{i\in B}x_i
        =\sum_{\varnothing\neq A\subseteq B}
        (-2)^{|A|-1}x_A.
    \end{equation}
    Modulo $2^k$, all terms with $|A|>k$ vanish.  Applying
    \cref{eq:parity-to-weighted} to each term in
    \cref{eq:phase-polynomial-sum-parities} therefore produces a weighted
    polynomial of the form \eqref{eq:weighted-phase-polynomial}.

    Conversely, for every nonempty $A\subseteq[n]$, Boolean
    inclusion--exclusion gives
    \begin{equation}
        \label{eq:weighted-to-parity}
        2^{|A|-1}x_A
        =\sum_{\varnothing\neq B\subseteq A}
        (-1)^{|B|-1}
        \bigoplus_{i\in B}x_i.
    \end{equation}
    Each parity on the right is
    $\langle \sum_{i\in B}e_i,x\rangle$, where $e_i$ is the $i$th
    standard basis vector.  Substituting \cref{eq:weighted-to-parity} into
    \cref{eq:weighted-phase-polynomial} yields
    \cref{eq:phase-polynomial-sum-parities}.
\end{proof}

We call the equivalent functions in
\Cref{thm:phase-polynomial-equivalence} \emph{$k$th-order phase
polynomials}.

The parity-sum form makes closure under changes of binary coordinates
immediate.

\begin{corollary}[Affine pullback]
    \label{cor:phase-polynomial-affine-pullback}
    If $q:\F_2^n\to\Z_{2^k}$ is a $k$th-order phase polynomial and
    $\phi(y)=Ay\oplus b$ is an affine map from \(\F_2^m\) to \(\F_2^n\),
    then $q\circ\phi$ is a $k$th-order phase polynomial on \(\F_2^m\).
\end{corollary}

\begin{proof}
    For $\beta\in\F_2$ and $t\in\{0,1\}$,
    \(\beta\oplus t=\beta+(1-2\beta)t\).  Hence
    \[
        \langle u,Ay\oplus b\rangle
        =\langle u,b\rangle
        +\bigl(1-2\langle u,b\rangle\bigr)
        \langle A^{\mathsf T}u,y\rangle.
    \]
    Substitution into the parity-sum form of
    \Cref{thm:phase-polynomial-equivalence} gives another constant plus a
    sum of parity functions.
\end{proof}

We next record how phase polynomials behave under discrete differentiation.
For $v\in\F_2^n$, recall that
\[
    \Delta_vq(x)=q(x\oplus v)-q(x).
\]

\begin{theorem}[Discrete derivatives of phase polynomials]
    \label{thm:phase-polynomial-derivative}
    Let $q:\F_2^n\to\Z_{2^k}$ be a $k$th-order phase polynomial, with
    $k\geq2$.  For every $v\in\F_2^n$, there are a constant
    $r_v\in\Z_{2^k}$ and a $(k-1)$st-order phase polynomial
    $q_v:\F_2^n\to\Z_{2^{k-1}}$ such that
    \begin{equation}
        \label{eq:general-phase-derivative}
        \Delta_vq(x)=r_v+2q_v(x)\pmod {2^k}.
    \end{equation}
    More generally, an $s$-fold discrete derivative, for $1\leq s<k$,
    is a constant plus $2^s$ times a $(k-s)$th-order phase polynomial;
    every $k$-fold derivative is constant.
\end{theorem}

\begin{proof}
    Write $q$ in the parity-sum form
    \eqref{eq:phase-polynomial-sum-parities}.  Since
    \[
        \langle u,x\oplus v\rangle-\langle u,x\rangle
        =\begin{cases}
            0, & \langle u,v\rangle=0,\\
            1-2\langle u,x\rangle, & \langle u,v\rangle=1,
        \end{cases}
    \]
    we obtain
    \begin{equation}
        \label{eq:derivative-parity-form}
        \Delta_vq(x)
        =\sum_{\langle u,v\rangle=1}a_u
        -2\sum_{\langle u,v\rangle=1}a_u\langle u,x\rangle
        \pmod {2^k}.
    \end{equation}
    The first sum is constant.  After reducing its coefficients modulo
    $2^{k-1}$, the second sum is a phase polynomial over
    \(\Z_{2^{k-1}}\), proving \cref{eq:general-phase-derivative}.
    Iterating this argument proves the remaining statements.
\end{proof}

For $k=3$, a discrete derivative is therefore a Clifford phase polynomial.
We make the corresponding normal forms explicit.

\begin{theorem}[Clifford phase polynomials]
    \label{thm:clifford-phase-polynomial}
    A function $d:\F_2^n\to\Z_8$ is the phase function of a diagonal
    Clifford operator, up to global phase, if and only if it has either of
    the equivalent forms
    \begin{align}
        d(x)
        &=d_0+2\sum_i \alpha_i x_i
          +4\sum_{i<j}\beta_{ij}x_ix_j \pmod 8,
          \label{eq:clifford-phase-weighted}\\
        &=d_0+2\sum_{u\in\F_2^n\setminus\{0\}}
          b_u\langle u,x\rangle \pmod 8,
          \label{eq:clifford-phase-parity}
    \end{align}
    where \(d_0\in\Z_8\), \(\alpha_i,b_u\in\Z_4\), and
    \(\beta_{ij}\in\F_2\).
\end{theorem}

\begin{proof}
    The diagonal Clifford group is generated, up to global phase, by the
    $\S_i$ and \(\CZ_{ij}\) gates.  Their phase functions are $2x_i$ and
    $4x_ix_j$, respectively, which proves
    \eqref{eq:clifford-phase-weighted}.  Conversely, every polynomial of
    that form is implemented by
    \(\prod_i \S_i^{\alpha_i}\prod_{i<j}\CZ_{ij}^{\beta_{ij}}\).

    Finally, dividing the nonconstant part of
    \eqref{eq:clifford-phase-weighted} by $2$ gives a second-order phase
    polynomial over \(\Z_4\).  Applying
    \Cref{thm:phase-polynomial-equivalence} with $k=2$ proves its
    equivalence with \eqref{eq:clifford-phase-parity}.
\end{proof}

\begin{theorem}[Discrete derivatives of Clifford phase polynomials]
    \label{thm:clifford-phase-derivative}
    Let $d:\F_2^n\to\Z_8$ be a Clifford phase polynomial.  For every
    $v\in\F_2^n$, there are $r_v\in\Z_4$ and a linear function
    \(\ell_v:\F_2^n\to\F_2\) such that
    \begin{equation}
        \label{eq:clifford-phase-derivative}
        \Delta_vd(x)=2r_v+4\ell_v(x)\pmod 8.
    \end{equation}
\end{theorem}

\begin{proof}
    Use \cref{eq:clifford-phase-parity} and write
    $d(x)=d_0+2\sum_u b_u\langle u,x\rangle$.
    \cref{eq:derivative-parity-form} gives
    \[
        \Delta_vd(x)
        =2\sum_{\langle u,v\rangle=1}b_u
        -4\sum_{\langle u,v\rangle=1}b_u\langle u,x\rangle
        \pmod 8.
    \]
    Take $r_v=\sum_{\langle u,v\rangle=1}b_u\pmod4$ and
    \[
        \ell_v(x)=
        \bigoplus_{\substack{u:\,\langle u,v\rangle=1\\ b_u\text{ odd}}}
        \langle u,x\rangle.
    \]
    The minus sign is immaterial after multiplication by $4$ modulo $8$,
    and \(\ell_v\) is linear.
\end{proof}
